\documentclass[11pt]{article}

\usepackage[margin=0.9in]{geometry}
\usepackage{parskip}
\usepackage{titlesec}
\usepackage{enumitem}
\usepackage{xcolor}
\usepackage{hyperref}

\hypersetup{
    colorlinks=true,
    linkcolor=blue!60!black,
    urlcolor=blue!60!black,
    citecolor=blue!60!black
}

\titleformat{\section}{\normalsize\bfseries}{\thesection.}{0.5em}{}
\titlespacing*{\section}{0pt}{0.8em}{0.3em}

\setlist[itemize]{leftmargin=1.3em, itemsep=1pt, topsep=1pt}
\setlist[enumerate]{leftmargin=1.5em, itemsep=1pt, topsep=1pt}

\title{\vspace{-2em}\textbf{ESE 545 Final Project Proposal}\\[0.2em]
       \large Hypothesis-Driven Code Debugging Agent}
\author{[NAME 1 / ID] \quad [NAME 2 / ID] \quad (2-person team)}
\date{April 12, 2026}

\begin{document}
\maketitle
\vspace{-2em}

%====================================================================
\section{Problem Description}

I propose to design and implement an \textbf{LLM-based code debugging agent} that, given a Python file and a failing test, iteratively localizes the bug and produces a patch that makes the tests pass. The agent is not a single LLM call that ``rewrites the buggy function.'' It is an explicit control loop in which the LLM serves as a reasoning engine---generating hypotheses, proposing patches, and reading code---while a \textbf{symbolic planner} outside the LLM manages the investigation: which hypothesis to pursue, which tool to call, when to commit a patch, when to back out, and when to stop.

The motivation is that modern LLMs are fluent code generators but weak at \emph{ruling things out}. Asked to ``fix this bug,'' an LLM will often produce a plausible-looking patch that doesn't address the real cause, and when shown it still fails, will propose another plausible-looking patch that repeats the mistake. The agent introduces \textbf{memory, scoring, and a discipline of elimination} at the planner level, so that wrong hypotheses get demoted and the search space actually shrinks over time. The LLM proposes; the test---not the LLM---judges.

\textbf{Scope.} The agent targets localized, reproducible bugs in a single Python file where a failing test provides a ground-truth signal: off-by-one errors, wrong operators or constants, missing edge cases, incorrect conditions, wrong variable references, and simple type errors. Multi-file bugs, concurrency, performance, flaky tests, and environment issues are out of scope.

%====================================================================
\section{Implementation}

\textbf{Components.} The agent consists of:
\begin{itemize}
    \item \textbf{State / Memory.} A structured record holding the target file, test command, current traceback, a hypothesis queue (each with a confidence score and status), patches already tried, and a full action/observation history.
    \item \textbf{Skills (tools).} Pure Python functions: \texttt{run\_tests}, \texttt{read\_file}, \texttt{search\_code}, \texttt{get\_traceback\_frames}, \texttt{run\_snippet} (executes a probe to check a suspected variable state), \texttt{apply\_patch}, \texttt{revert\_patch}, \texttt{lint}. Each returns a structured observation.
    \item \textbf{LLM interface.} Two structured-output calls: \texttt{generate\_hypotheses(state)} and \texttt{propose\_patch(state, hypothesis)}, plus \texttt{pick\_investigation} which chooses a probe tool from a constrained menu. The LLM cannot invent new tools or control the loop.
    \item \textbf{Planner (symbolic).} A deterministic Python function \texttt{plan\_next\_action(state)} that inspects the hypothesis queue and decides the next action. No LLM call is made inside the planner.
    \item \textbf{Execution loop.} Receives the action from the planner, invokes the appropriate LLM or tool call, updates state (including hypothesis scores), and repeats until tests pass, hypotheses are exhausted, or an iteration cap is reached.
\end{itemize}

\textbf{Inputs / Outputs.} Inputs: a path to the buggy Python file and a shell command running a failing test (typically \texttt{pytest}), plus optional iteration cap and confidence threshold. Outputs: (1) a patched file on disk, (2) a structured trace log of all hypotheses, tool calls, and observations, and (3) a final status (fixed, exhausted, or capped).

\textbf{Execution cycle.}
\begin{enumerate}
    \item If no hypotheses exist $\rightarrow$ call LLM to generate a ranked list of candidate causes from the traceback and surrounding code.
    \item Planner selects the top \emph{open} hypothesis.
    \item If its score exceeds the confidence threshold $\rightarrow$ propose a patch, apply it, run tests. If tests pass, done. If not, revert and penalize.
    \item Otherwise $\rightarrow$ investigate it with a tool call (e.g.\ \texttt{run\_snippet}). Update the score based on the observation.
    \item If all hypotheses are exhausted, regenerate with the ruled-out list as negative context. If this recurs, terminate.
\end{enumerate}

\textbf{Hypothesis scoring.} Each hypothesis carries a numeric confidence updated by fixed rules---e.g.\ $+2$ if a probe confirmed the suspected bad state, $-3$ if a probe showed the value was fine, $-5$ if a patch attempt failed, $-2$ if a patch changed the failing signal but did not fix it. This gives the agent the ability to eliminate candidates, which a stateless LLM call lacks.

\textbf{Action selection rules.} \texttt{plan\_next\_action} is deterministic Python:
\begin{itemize}
    \item No hypotheses $\rightarrow$ generate.
    \item Top hypothesis above threshold $\rightarrow$ patch.
    \item Top hypothesis open but below threshold $\rightarrow$ investigate.
    \item All hypotheses closed $\rightarrow$ regenerate with exclusions, or give up past a regeneration cap.
\end{itemize}
These rules live in code, not in a prompt---the key distinction from a ReAct-style pipeline where the LLM itself chooses the next step.

\textbf{Backtracking.} Every patch is applied through \texttt{apply\_patch} and tracked with a patch ID. A failed patch is reverted via \texttt{revert\_patch} before the loop continues, making the search over patch space safe and reversible.

\textbf{Where planning enters.} Planning enters at three points: (1) \emph{action selection}---the planner, not the LLM, decides whether the next step is generate, investigate, patch, or stop; (2) \emph{hypothesis lifecycle}---scoring and status transitions (\texttt{open} $\rightarrow$ \texttt{failed} / \texttt{confirmed} / \texttt{ruled\_out}) are symbolic state updated by deterministic rules; (3) \emph{search control}---the planner enforces the iteration cap, regeneration cap, and backtracking policy, and the LLM has no ability to bypass them. The LLM is confined to what it does well (reading code, generating hypotheses, proposing patches) while the planner handles what LLMs do poorly: tracking what's been tried, resisting repeated failed patches, and knowing when to stop.

\end{document}
